
\uline{by} @@M0@@. \uline{If} @@M1@@ \uline{is the structural morphism of} @@M2@@, \uline{then} @@M3@@, \uline{defined because} @@M4@@ \uline{is flat}, \uline{is surjective and defines an isomorphism}:
@@M5@@

Let @@M6@@, let @@M7@@ be the evident basis of @@M8@@, and let @@M9@@ be the sub-Module of @@M10@@ generated by @@M11@@. The subfunctor of @@M12@@ that assigns to every @@M13@@-scheme @@M14@@ the subgroup of automorphisms of @@M15@@ leaving invariant the submodules @@M16@@ for every @@M17@@ is representable by a group subscheme @@M18@@ of @@M19@@, which is a Borel subgroup of @@M20@@. The quotient sheaf @@M21@@ (for the fpqc topology) is representable by the scheme of flags of @@M22@@, which will be denoted by @@M23@@. Let @@M24@@ be the canonical quotients of the inverse image of @@M25@@ on @@M26@@, let @@M27@@ be the kernel of @@M28@@, and let @@M29@@.

Moreover, @@M30@@ is the semidirect product of its unipotent radical @@M31@@ and its maximal subtorus @@M32@@ leaving the @@M33@@ invariant. Let @@M34@@. Then @@M35@@ is a torsor with base @@M36@@ under the group @@M37@@. If @@M38@@ is an @@M39@@-scheme, giving an @@M40@@-point of @@M41@@ is equivalent to giving, over an @@M42@@ faithfully flat and quasi-compact over @@M43@@, an automorphism of @@M44@@ modulo the unipotents of @@M45@@, that is, a flag of @@M46@@ whose submodules are free and, for every @@M47@@, a class of elements @@M48@@ free over @@M49@@ such that @@M50@@; equivalently, it is giving a flag of @@M51@@ and, for every @@M52@@, a basis of @@M53@@, these data being subject to the gluing condition on @@M54@@; by fpqc descent, this amounts to the same data on @@M55@@. One then sees immediately that @@M56@@ is the affine @@M57@@-scheme defined by the algebra:

